
En effet, dans (i), la premi\`ere assertion r\'esulte de 6.7 (i) ; la deuxi\`eme, de 6.10 et 1.8 (i), et de 6.2 (appliqu\'e \`a @@M0@@, @@M1@@) et 1.3. Enfin, (ii) r\'esulte de (i) et 1.13.

\textbf{Remarque 6.12.} Soient @@M2@@ un corps alg\'ebriquement clos, et @@M3@@ un @@M4@@-sch\'ema projectif purement de dimension @@M5@@, qui est sans cycle premier associ\'e immerg\'e. Si l'on admet des \'el\'ements nilpotents dans @@M6@@ de fa\c{c}on arbitraire, on ne peut plus borner les coefficients dans le polyn\^ome de Hilbert @@M7@@ en termes de @@M8@@ seulement.

Par exemple, soit @@M9@@ une courbe lisse de degr\'e @@M10@@ dans @@M11@@. Pour tout @@M12@@, il existe une surface lisse @@M13@@ de degr\'e @@M14@@, qui contient @@M15@@. Alors @@M16@@ est un diviseur de Cartier sur @@M17@@; soit @@M18@@. Enfin, on a la suite exacte @@M19@@ et la formule @@M20@@ ; d'o\`u que @@M21@@ (@@M22@@), bien que @@M23@@.

\textbf{Corollaire 6.13.} \emph{Sous les conditions de 6.1, consid\'erons une suite exacte}
@@M24@@
\emph{avec @@M25@@. Alors pour @@M26@@,}
@@M27@@

En effet, soit @@M28@@ ; \'evidemment, @@M29@@ est un @@M30@@-faisceau, et @@M31@@ ; d'o\`u l'assertion gr\^ace \`a 6.7 (i), et la formule @@M32@@.

